\section*{Exercises on \S\arabic{exercisegroup}}
\phantomsection
\addcontentsline{toc}{section}{Exercises on \protect\S\arabic{exercisegroup}}

\begin{enumerate}
\def\labelenumi{\arabic{enumi})}
\item
  Let \(f\) be a real function defined and left continuous on an open interval \(I = ]a, b[\) in \(R\) ; suppose that at all the points of the complement \(B\) with respect to \(I\) of a countable subset of \(I\) the function \(f\) is increasing to the right, that is, at every point \(x \in B\) there exists a \(y\) such that \(x < y \leqslant b\) and such that for all \(z\) such that \(x \leqslant z < y\) one has \(f(x) \leqslant f(z)\) . Show that \(f\) is increasing on \(I\) (argue as in prop. 2).
\item
  In the field \(Q_{p}\) of p-adic numbers (Gen.~Top., III, p.~322, exerc. 23) every p-adic integer \(x \in Z_{p}\) has one and only one expansion in the form \(x = a_{0} + a_{1}p + \cdots + a_{n}p^{n} + \cdots\) , where the \(a_{j}\) are rational integers such that \(0 \leqslant a_{j} \leqslant p - 1\) for each j. For each \(z \in Z_{p}\) put
\end{enumerate}

\[
f (x) = a _ {0} + a _ {1} p ^ {2} + \dots + a _ {n} p ^ {2 n} + \dots ;
\]

show that, on \(Z_{p}\) , f is a continuous function which is not constant on a neighbourhood of any point yet has a zero derivative at every point.

\begin{enumerate}
\def\labelenumi{\arabic{enumi})}
\setcounter{enumi}{2}
\tightlist
\item
  \begin{enumerate}
  \def\labelenumii{\alph{enumii})}
  \tightlist
  \item
    Let K be the triadic Cantor set (Gen.~Top., IV, p.~338), let \(I_{n,p}\) be the \(2^{n}\) contiguous intervals of K with length \(1/3^{n+1}\) ( \(1 \leqslant p \leqslant 2^{n}\) ), and \(K_{n,p}\) the \(2^{n+1}\) closed intervals of length \(1/3^{n+1}\) whose union is the complement of the union of the \(I_{m,p}\) for \(m \leqslant n\) . Let \(\alpha\) be a number such that \(1 < \alpha < 3/2\) ; for each n we denote by \(f_{n}\) the continuous increasing function on {[}0, 1{]} which is equal to 0 for x = 0, constant on each of the intervals \(I_{m,p}\) for \(m \leqslant n\) , is affine linear on each of the intervals \(K_{n,p}\) ( \(1 \leqslant p \leqslant 2^{n+1}\) ) and such that \(f_{d}'(x) = \alpha^{n+1}\) on each of the interiors of these last intervals. Show that the series with general term \(f_{n}\) is uniformly convergent on {[}0, 1{]}, that its sum is a function f which admits a right derivative (finite or not) everywhere in {[}0, 1{[}, and that one has \(f_{r}'(x) = +\infty\) at every point of K distinct from the left-hand endpoints of the contiguous intervals \(I_{n,p}\) .
  \end{enumerate}
\end{enumerate}

\begin{enumerate}
\def\labelenumi{\alph{enumi})}
\setcounter{enumi}{1}
\tightlist
\item
  Let g be a continuous increasing map of \([0, 1]\) onto itself, constant on each of the intervals \(I_{n,p}\) (Gen.~Top., IV, p.~403, exerc. 9). If \(h = f + g\) , show that h admits a right derivative equal to \(f_{d}^{\prime}(x)\) at every point x of \([0, 1]\) .
\end{enumerate}

\begin{enumerate}
\def\labelenumi{\arabic{enumi})}
\setcounter{enumi}{3}
\tightlist
\item
  Let \(f\) be a finite real function, continuous on a compact interval \([a, b]\) in \(\mathbf{R}\) , and having a right derivative at every point of the open interval \(]a, b[\) . Let \(m\) and \(M\) be the greatest lower bound and least upper bound (finite or not) of \(f_d'\) over \(]a, b[\) .
\end{enumerate}

\begin{enumerate}
\def\labelenumi{\alph{enumi})}
\item
  Show that when x and y run through {]}a, b{[} keeping \(x \neq y\) , the set of values of \((f(x) - f(y))/(x - y)\) contains {]}m, M{[} and is contained in {[}m, M{]}. (Reduce to proving that if \(f_{d}^{\prime}\) takes two values of opposite sign at the two points c, d of {]}a, b{[} (with c \textless{} d), then there exist two distinct points of the interval {]}c, d{[} where f takes the same value).
\item
  If, further, f has a left derivative at every point of \(]a\) , b{[} then the infima (resp. suprema) of \(f_{d}^{\prime}\) and \(f_{g}^{\prime}\) over \(]a\) , b{[} are equal.
\item
  Deduce that if \(f\) is differentiable on \(]a, b[\) then the image under \(f'\) of every interval contained in \(]a, b[\) is itself an interval, and consequently connected (use \(a\) ).
\end{enumerate}

\begin{enumerate}
\def\labelenumi{\arabic{enumi})}
\setcounter{enumi}{4}
\item
  Let \(\mathbf{f}\) be the vector mapping of \(I = [0, 1]\) into \(\mathbf{R}^3\) defined as follows: for \(0 \leqslant t \leqslant \frac{1}{4}\) , \(\mathbf{f}(t) = (-4t, 0, 0)\) ; for \(\frac{1}{4} \leqslant t \leqslant \frac{1}{2}\) let \(\mathbf{f}(t) = (-1, 4t - 1, 0)\) ; for \(\frac{1}{2} \leqslant t \leqslant \frac{3}{4}\) let \(\mathbf{f}(t) = (-1, 1, 4t - 2)\) ; finally, for \(\frac{3}{4} \leqslant t \leqslant 1\) take \(\mathbf{f}(t) = (4t - 4, 1, 1)\) . Show that the convex set generated by the set \(\mathbf{f}_d'(\mathrm{I})\) is not identical to the closure of the set of values of \(\frac{\mathbf{f}(y) - \mathbf{f}(x)}{y - x}\) as \((x, y)\) runs through the set of pairs of distinct points of \(I\) (cf.~exerc. \(4a\) )).
\item
  On the interval \(\mathbf{I} = [-1, + 1]\) consider the vector function \(\mathbf{f}\) , with values in \(\mathbf{R}^2\) , defined as follows: \(\mathbf{f}(t) = (0,0)\) for \(-1\leqslant t\leqslant 0\) ;
\end{enumerate}

\[
\mathbf {f} (t) = \left(t ^ {2} \sin \frac {1}{t}, t ^ {2} \cos \frac {1}{t}\right)
\]

for \(0 \leqslant t \leqslant 1\) . Show that \(\mathbf{f}\) is differentiable on \([-1, +1[\) but that the image of this interval under \(\mathbf{f}'\) is not a connected set in \(\mathbf{R}^2\) ( \(cf.\) exerc. \(4c\) )).

\begin{enumerate}
\def\labelenumi{\arabic{enumi})}
\setcounter{enumi}{6}
\item
  Let f be a continuous vector function defined on an open interval \(I \subset R\) , with values in a normed space E over R, and admitting a right derivative at every point of I. Show that the set of points of I where f admits a derivative is the complement of a meagre subset of I (use exerc. 5 b) of I, p.~36, and prop. 6 of I, p.~18).
\item
  Consider, on the interval {[}0, 1{]}, a family \((\mathrm{I}_{n,p})\) of pairwise disjoint open intervals, defined inductively as follows: the integer \(n\) takes all values \(\geqslant 0\) ; for each value of \(n\) the integer \(p\) takes the values \(1,2,\ldots ,2^n\) ; one has \(\mathrm{I}_{0,1} = ]\frac{1}{3},\frac{2}{3}[\) ; if \(\mathrm{J}_n\) is the union of the intervals \(\mathrm{I}_{m,p}\) corresponding to the numbers \(m\leqslant n\) , the complement of \(\mathrm{J}_n\) is the union of \(2^{n + 1}\) pairwise disjoint closed intervals \(\mathrm{K}_{n,p}\) ( \(1\leqslant p\leqslant 2^{n + 1}\) ). If \(\mathrm{K}_{n,p}\) is an interval \([a,b]\) one then takes for \(\mathrm{I}_{n + 1,p}\) the open interval with endpoints \(b - \frac{b - a}{3}\left(1 + \frac{1}{2^n}\right)\) and \(b - \frac{b - a}{3.2^n}\) . Let E be the perfect set which is the complement with respect to {[}0, 1{]} of the union of the \(\mathrm{I}_{n,p}\) . Define on {[}0, 1{]} a continuous real function \(f\) which admits a right derivative at every point of {[}0, 1{]}, but fails to have a left derivative at the uncountable subset of E of points distinct from the endpoints of intervals contiguous to E (cf.~exerc. 7). (Take \(f(x) = 0\) on E, define \(f\) suitably on each of the intervals \(\mathrm{I}_{n,p}\) in such a way that for every \(x\in \mathrm{E}\) there are points \(y < x\) not belonging to E, arbitrarily close to \(x\) , and such that \(\frac{f(y) - f(x)}{y - x} = -1\) .)
\item
  Let \(f\) and \(g\) be two finite real functions, continuous on \([a, b]\) , both having a finite derivative on \(]a, b[\) ; show that there exists a \(c\) such that \(a < c < b\) and that
\end{enumerate}

\[
\left| \begin{array}{c c} f (b) - f (a) & g (b) - g (a) \\ f ^ {\prime} (c) & g ^ {\prime} (c) \end{array} \right| = 0.
\]

¶ 10) Let \(f\) and \(g\) be two finite real functions, strictly positive, continuous and differentiable on an open interval I. Show that if \(f'\) and \(g'\) are strictly positive and \(f'/g'\) is strictly increasing on I, then either \(f/g\) is strictly increasing on I, or else there exists a number \(c \in I\) such that \(f / g\) is strictly decreasing for \(x \leqslant c\) and strictly increasing for \(x \geqslant c\) (note that if one has \(f'(x) / g'(x) < f(x) / g(x)\) then also

\[
f ^ {\prime} (y) / g ^ {\prime} (y) <   f (y) / g (y)
\]

for all y \textless{} x).

\begin{enumerate}
\def\labelenumi{\arabic{enumi})}
\setcounter{enumi}{10}
\tightlist
\item
  Let \(f\) be a complex function, continuous on an open interval \(I\) , vanishing nowhere, and admitting a right derivative at every point of \(I\) . For \(|f|\) to be increasing on \(I\) it is necessary and sufficient that \(\mathcal{R}(f_d' / f) \geqslant 0\) on \(I\) .
\end{enumerate}

¶ 12) Let \(f\) be a differentiable real function on an open interval I, \(g\) its derivative on I, and \([a, b]\) a compact interval contained in I; suppose that \(g\) is differentiable on the open interval \([a, b]\) but not necessarily right (resp. left) continuous at the point \(a\) (resp. \(b\) ); show that there exists \(c\) such that \(a < c < b\) and that

\[
g (b) - g (a) = (b - a) g ^ {\prime} (c)
\]

(use exerc. 4 c) of I, p.~36).

\begin{enumerate}
\def\labelenumi{\arabic{enumi})}
\setcounter{enumi}{12}
\tightlist
\item
  One terms the symmetric derivative of a vector function \(\mathbf{f}\) at a point \(x_0\) interior to the interval where \(\mathbf{f}\) is defined, the limit (when it exists) of \(\frac{\mathbf{f}(x_0 + h) - \mathbf{f}(x_0 - h)}{2h}\) as \(h\) tends to 0 remaining \(>0\) .
\end{enumerate}

\begin{enumerate}
\def\labelenumi{\alph{enumi})}
\item
  Generalize to the symmetric derivative the rules of calculus established in § 1 for the derivative.
\item
  Show that theorems 1 and 2 of § 2 remain valid when one replaces the words ``right derivative'' by ``symmetric derivative''.
\end{enumerate}

\begin{enumerate}
\def\labelenumi{\arabic{enumi})}
\setcounter{enumi}{13}
\tightlist
\item
  Let \(\mathbf{f}\) be a vector function defined and continuous on a compact interval \(\mathrm{I} = [a, b]\) in \(\mathbf{R}\) , with values in a normed space over \(\mathbf{R}\) . Suppose that \(\mathbf{f}\) admits a right derivative at all points of the complement with respect to \([a, b[\) of a countable subset A of this interval. Show that there exists a point \(x \in ]a\) , \(b[\cap ∁A\) such that
\end{enumerate}

\[
\left\| \mathbf {f} (b) - \mathbf {f} (a) \right\| \leqslant \left\| \mathbf {f} _ {d} ^ {\prime} (x) \right\| (b - a).
\]

(Argue by contradiction, decomposing \([a, b[\) into three intervals \([a, t[, [t, t + h[\) and \([t + h, b[\) with \(t \notin A\) ; if \(k = \| \mathbf{f}(b) - \mathbf{f}(a)\| / (b - a)\) , note that for \(h\) sufficiently small one has \(\| \mathbf{f}(t + h) - \mathbf{f}(t)\| < k.h\) , and use th. 2 of I, p.~15 for the other intervals.)

\setcounter{exercisegroup}{2}
\refstepcounter{exercisegroup}
